%% Örnek Özgeçmiş - [YOUR_NAME]
%% Tüm yetkinlikleri, deneyimleri ve başarıları kapsayan ana özgeçmiş.
%% Pozisyona özel özgeçmişler hazırlanırken referans olarak kullanılır.
%%
%% KURULUM: Kendi bilgilerinle doldurmak için /setup komutunu çalıştır.
%% Derleme: lualatex main_example.tex
%% (pdflatex, yeni MiKTeX kurulumlarında fontawesome5 hatası verebilir.)
%%
%% Not: İngilizce ilanlara başvururken /apply bu şablonu İngilizceye çevirip
%% bölüm başlıklarını da İngilizce yazar (Experience, Education vb.).

\documentclass[11pt,a4paper,sans]{moderncv}
\moderncvstyle{banking}
\moderncvcolor{blue}

% İsim ve bölüm başlıklarını moderncv mavisiyle (color1) göster. banking
% stili lualatex ile bunları siyah bırakabiliyor.
\renewcommand*{\firstnamestyle}[1]{{\fontsize{34}{36}\bfseries\upshape\color{color1}#1}}
\renewcommand*{\lastnamestyle}[1]{{\fontsize{34}{36}\bfseries\upshape\color{color1}#1}}
\renewcommand*{\sectionstyle}[1]{{\sectionfont\color{color1}#1}}

\usepackage[utf8]{inputenc}
\usepackage{hyperref}
\hypersetup{
    colorlinks=true,
    linkcolor=blue,
    filecolor=magenta,
    urlcolor=blue,
    pdftitle={[YOUR_NAME] - Özgeçmiş},
    pdfpagemode=FullScreen,
}
\usepackage[scale=0.80]{geometry}
\usepackage{import}

% kişisel bilgiler
\name{[Ad]}{[Soyad]}
\address{[İlçe, İl]}{}{}
\phone[mobile]{[+90 5XX XXX XX XX]}
\email{[e-posta@ornek.com]}
\extrainfo{\href{[https://linkedin.com/in/profilin]}{LinkedIn}, \href{[https://github.com/kullanici-adin]}{GitHub}}

\begin{document}

\makecvtitle

% ============================================================
%     PROFİL
% ============================================================

\vspace{6pt}
\small{[Başvurduğun pozisyona göre 3-5 satırlık bir profil yazısı. Seni bu iş için öne çıkaran şeye odaklan. Örnek: "[Alan] bölümü mezunu, [X] yıllık deneyime sahip veri analisti. [Sektör] bilgisini Python ve SQL ile uygulamalı analiz becerisiyle birleştiriyor. Verinin toplanmasından yöneticilere sunulan raporlara kadar uçtan uca süreçler kurdu."]}

% ============================================================
%     YETKİNLİKLER
% ============================================================

\section{Yetkinlikler}
\vspace{1pt}
\begin{itemize}

\item \textbf{[Yetkinlik Grubu 1]}: [Somut beceriler, araçlar, teknolojiler.]

\item \textbf{[Yetkinlik Grubu 2]}: [Somut beceriler, araçlar, teknolojiler.]

\item \textbf{[Yetkinlik Grubu 3]}: [Somut beceriler, araçlar, teknolojiler.]

\item \textbf{[Yetkinlik Grubu 4]}: [Sektör bilgisi, yöntemler, yaklaşımlar.]

\item \textbf{[Yetkinlik Grubu 5]}: [Programlar, platformlar, yazılımlar.]

\end{itemize}

% ============================================================
%     İŞ DENEYİMİ
% ============================================================

\section{İş Deneyimi}
\vspace{3pt}
\begin{itemize}

% --- Son pozisyon ---
\item{\cventry{[YYYY--Halen]}{[Pozisyon]}{[Şirket]}{[İl]}{}{\vspace{1pt}
\begin{itemize}
    \item [Başarı veya sorumluluk 1: somut ol, mümkünse rakam ver]
    \item [Başarı veya sorumluluk 2]
    \item [Başarı veya sorumluluk 3]
    \item [Başarı veya sorumluluk 4]
\end{itemize}}}

\vspace{3pt}

% --- Önceki pozisyon ---
\item{\cventry{[YYYY--YYYY]}{[Pozisyon]}{[Şirket]}{[İl]}{}{\vspace{1pt}
\begin{itemize}
    \item [Başarı veya sorumluluk 1]
    \item [Başarı veya sorumluluk 2]
    \item [Başarı veya sorumluluk 3]
\end{itemize}}}

\vspace{3pt}

% --- Daha eski pozisyon ---
\item{\cventry{[YYYY--YYYY]}{[Pozisyon]}{[Şirket]}{[İl]}{}{\vspace{1pt}
\begin{itemize}
    \item [Başarı veya sorumluluk 1]
    \item [Başarı veya sorumluluk 2]
\end{itemize}}}

\end{itemize}

% ============================================================
%     EĞİTİM
% ============================================================

\section{Eğitim}
\vspace{1pt}
\begin{itemize}

\item{\cventry{[YYYY--YYYY]}{[Bölüm], [Lisans / Yüksek Lisans]}{[Üniversite]}{[İl]}{}{\vspace{1pt}
Tez: ``[Tez Başlığı].'' [Araştırma konusunun kısa açıklaması.]
}}

\vspace{3pt}

\item{\cventry{[YYYY--YYYY]}{[Bölüm], [Lisans / Ön Lisans]}{[Üniversite]}{[İl]}{}{\vspace{1pt}
[Kısa açıklama veya öne çıkan dersler.]
}}

\end{itemize}

% ============================================================
%     YABANCI DİL
% ============================================================

\section{Yabancı Dil}
\vspace{1pt}
\begin{itemize}
\item Türkçe (ana dil), [Dil 2] ([seviye, örn. ileri / C1, YDS veya IELTS puanı]), [Dil 3] (orta).
\end{itemize}

% ============================================================
%     SERTİFİKALAR (isteğe bağlı)
% ============================================================

\section{Sertifikalar}
\vspace{1pt}
\begin{itemize}
\item{\textbf{[Sertifika Adı]} -- [Kurum] ([Yıl]).}
\end{itemize}

% ============================================================
%     ÖDÜLLER (isteğe bağlı)
% ============================================================

\section{Ödüller}
\vspace{1pt}
\begin{itemize}
\item{\textbf{[Ödül Adı]} -- [Etkinlik / Kurum] ([Yıl]).}
\end{itemize}

% ============================================================
%     EK BİLGİLER (isteğe bağlı, Türkiye'deki başvurularda yaygın)
%     Askerlik ve ehliyet çoğu yerli ilanda sorulur. İngilizce veya
%     yurt dışı başvurularda bu bölümü çıkar. Fotoğraf, medeni durum ve
%     doğum tarihi zorunlu değildir; eklemek tamamen senin tercihin.
% ============================================================

\section{Ek Bilgiler}
\vspace{1pt}
\begin{itemize}
\item \textbf{Askerlik}: [Tamamlandı / Tecilli / Muaf]
\item \textbf{Ehliyet}: [B sınıfı]
\end{itemize}

% ============================================================
%     REFERANSLAR
% ============================================================

\section{Referanslar}
\vspace{1pt}
\begin{itemize}
\item{Talep edilmesi hâlinde sunulacaktır.}
\end{itemize}

\end{document}
